The idea that the universe is expanding is of recent origin. 
All the early cosmologies were essentially stationary and even Einstein, whose theory of relativity is the basis for almost all modern developments in cosmology, found it natural to suggest a static model of the universe. 

However, there is a very grave difficulty associated with a static model such as Einstein's, which is supposed to have existed for an infinite time. 
For, if the stars had been radiating energy at their present rates for an infinite time, they would have needed an infinite supply of energy. 
Further, the flux of radiation now would be infinite. 

Alternatively, if they had only a limited supply of energy, the whole universe would by now have reached thermal equilibrium, which is certainly not the case. 
This difficulty was noticed by Olbers, who however was not able to suggest any solution. 
The discovery of the recession of the nebulae by Hubble led to the abandonment of static models in favour of ones which were expanding. 

Clearly there are several possibilities: the universe may have expanded from a highly dense state a finite time ago (the so-called ``big bang'' model); another is that the present expansion may have been preceded by a contraction which, in Its turn may have been preceded by another expansion (the ``bouncing'' or oscillating model); however, this model suffers from the same difficulties over entropy as the static model. Finally, it is possible that the expansion may have been proceeding at much the same rate for an infinite time. It is then necessary to postulate some form of continual creation of matter in order to prevent the expansion from reducing the density to zero. This leads to the ``steady-state'' model which, although expanding, presents the same appearance at all times.

The early cosmologies naturally placed man at or near the centre of the universe, but since the time of Copernicus we have been demoted to a medium-sized planet going round a medium-sized star somewhere near the edge of a fairly average galaxy. We are now so humble that we would not claim to occupy any special position. However, observations seem to indicate that, within experimental error (which is fairly high), galaxies have a spatially isotropic distribution around us. As we are not claiming any special position, the distribution must be isotropic about every point. This implies that the distribution must be spatially homogeneous as well as isotropic. Of course, this homogeneity and isotropy hold only on a large scale; locally there are considerable departures from both.\\
Robertson and Walker have shown that the metric of a model 
that is spatially homogeneous and isotropic may be written in the form:

\[
ds^{2} = dt^{2} - R^{2}(t) \left[ \frac{dr^{2}}{1 - \kappa r^{2}} + r^{2} \left( d\theta^{2} + \sin^{2}\theta \, d\phi^{2} \right) \right],
\qquad \kappa = 0, \pm 1
\]

This form will be used extensively in the following chapters. 
In Chapter~1 it will be shown that the Hoyle--Narlikar theory 
of gravitation is incompatible with a metric of this form. 
In Chapter~2 perturbations of this form will be considered in 
a linearized approximation and, in Chapter~3, gravitational 
radiation will be considered in a model which tends asymptotically 
to this form.

Certain of the Robertson--Walker models possess ``horizons''. 
There are two types: particle horizons and event horizons. 
A particle horizon is said to exist when an observer's past 
light cone does not intersect the world line of every particle 
in the universe (or extended world line in the case of a 
particle which has been created). An example of a model with 
a particle horizon is the Einstein--de~Sitter model which has 

\[
\kappa = 0, \qquad R = t^{\tfrac{2}{3}} \, .
\]


This is a ``big--bang'' model as indeed are all the Robertson--Walker models that satisfy the Einstein equations:

\[
R_{ab} - \tfrac{1}{2} g_{ab} R = T_{ab}
\]

\noindent
\ldots and contain matter whose pressure is greater than minus one--third the density. 
An event horizon exists when there are events that a given observer will never see. 
The steady--state model $\bigl(\kappa = 0, \; R = e^{t}\bigr)$ is an example of one with an event horizon. 
Horizons will be further discussed in Chapter~4 which also deals with the occurrence of singularities of space--time and their connection with topology.

\medskip

Each chapter is self--contained and has its own references. 
The following notation is used throughout: space--time is taken to be a Riemannian manifold with metric tensor $g_{ij}$. 
This is taken to have signature $-2$ except in Chapter~2 where, in order to facilitate comparison with previous work, the signature is $+2$. 
Covariant differentiation is indicated by a semi--colon. 
Units are employed in which $c$, the speed of light, and $k$, the gravitational constant, equal one.
